\documentclass[11pt,a4paper]{article}
\usepackage[T1]{fontenc}
\usepackage[utf8]{inputenc}
\usepackage[english]{babel}
\usepackage{amsmath,amssymb}
\usepackage[margin=2cm]{geometry}
\setlength{\parindent}{0pt}
\setlength{\emergencystretch}{3em}
\allowdisplaybreaks

\title{About \texttt{index\_identities.tex}}
\author{}
\date{}
\begin{document}
\maketitle
This note explains what the test data in \texttt{index\_identities.tex} is used for. That file is input for the engine in \texttt{main.c}; it is read by the test \texttt{test\_index\_identities.c} in this folder.

\begin{enumerate}
\item \textbf{Purpose.} The file holds 77 true identities of tensor calculus written in index notation, in a fixed dimension: the Kronecker delta, the Levi-Civita symbol and components of arbitrary tensors, with the Einstein summation convention. The test asks the engine to prove each one, to write its proof, and to produce a certificate that the independent checker accepts.

\item \textbf{Where the statements come from.} Solved exercises: T.~Sochi, \emph{Solutions of Exercises of Principles of Tensor Calculus} (chapter 4, exercises 4.17, 4.21, 4.22, 4.26--4.30, 4.38, 4.39, 4.56, 4.70, 4.74, 4.76, 4.84; chapter 7, exercises 7.2--7.4, 7.8, 7.10), the index-notation handout of Ben-Yaacov and Roig, and problem set 1 of L.~Stein's Electromagnetism I (the double cross product, the Jacobi identity). Chains of equalities of the books are split into single identities, and a determinant is written out. Every statement was confirmed with a third computation that does not use the program: the oracle of \texttt{tools/fuzz/fuzz\_tensors.py}, run in its \texttt{data} mode on this file.

\item \textbf{What is being tested.} Statements made of $\delta$ and $\varepsilon$ alone, as $\varepsilon_{ijk}\varepsilon_{imn} = \delta_{jm}\delta_{kn} - \delta_{jn}\delta_{km}$, in two, three and four dimensions; and identities between arbitrary vectors and tensors, as the double cross product $\varepsilon_{ijk}\varepsilon_{klm} a_{j} b_{l} c_{m} = b_{i} a_{m} c_{m} - c_{i} a_{l} b_{l}$, where the components are independent unknowns.

\item \textbf{What the engine needs in order to pass.} To read which indices are free and which are summed, to find the dimension (the number of indices of $\varepsilon$, the quantifiers $\forall i, j \in \{1, 2, 3\}$, or a summation sign $\sum_{i=1}^{3}$), and to compare the two sides for every value of the free indices, as polynomials in the components.

\item \textbf{Format expected by the test.} One \texttt{\textbackslash item} per statement, with the formula written as \texttt{\$\textbackslash displaystyle ...\$}. Upper and lower indices are not distinguished; a number written alone as a side, as in $\varepsilon_{ijk} a_{j} a_{k} = 0$, is the zero tensor.

\item \textbf{How to run it.} From the top folder of the repository:
\begin{verbatim}
cc -std=c11 -O2 -Wall -Wextra -Werror -o /tmp/test_index_identities \
    tests/tensor_calculus/test_index_identities.c && /tmp/test_index_identities
\end{verbatim}

\item \textbf{Expected result.} The lines \texttt{PASS: 77 of 77 statements proved, each with a written proof}, \texttt{PASS: 77 of 77 statements proved, 77 with a certificate the checker accepts} and \texttt{PASS: 77 of 77 altered certificates rejected for the right reason}.
\end{enumerate}
\end{document}
